\begin{definition}[Pair frame]
\label{def:list-pair-frame}
For histories $u,t:\mathsf{BHist}$, define $\mathsf{PairFrame}(u,t)$ by
structural recursion on the first BEDC history:
$$
\begin{array}{rcl}
\mathsf{PairFrame}(\emp,t) &=& \Ezero(t),\\
\mathsf{PairFrame}(\Ezero(a),t) &=& \Eone(\Ezero(\mathsf{PairFrame}(a,t))),\\
\mathsf{PairFrame}(\Eone(a),t) &=& \Eone(\Eone(\mathsf{PairFrame}(a,t))).
\end{array}
$$
The outer $\Eone$ marks a nonempty first component, while the inner
$\Ezero/\Eone$ tag records the constructor of that first component.
\end{definition}
\leandef{BEDC.Derived.ListUp.PairFrame}

\begin{lemma}[Pair frame injectivity]
\label{lem:list-pair-frame-injectivity}
For histories $u,u',t,t':\mathsf{BHist}$, if
$$
\hsame(\mathsf{PairFrame}(u,t),\mathsf{PairFrame}(u',t')),
$$
then $\hsame(u,u')\land\hsame(t,t')$.
\end{lemma}
\begin{proof}
Proceed by structural induction and case analysis on $u$ and $u'$. The mixed
$\emp/\Ezero$ and $\emp/\Eone$ cases contradict BEDC constructor no-confusion
because $\mathsf{PairFrame}(\emp,t)$ begins with $\Ezero$, while the nonempty
cases begin with $\Eone$. The mixed $\Ezero/\Eone$ cases contradict
no-confusion after inverting the outer $\Eone$ and comparing the inner tags. In
the matching $\emp/\emp$ case, constructor inversion gives $\hsame(t,t')$. In
the matching $\Ezero/\Ezero$ and $\Eone/\Eone$ cases, invert the outer
$\Eone$, then the matching inner tag, and apply the induction hypothesis to the
recursive pair-frame equality. Combining the exposed recursive head equality
with the constructor equality gives $\hsame(u,u')$, and the induction
hypothesis also gives $\hsame(t,t')$.
\end{proof}
\leanchecked{BEDC.Derived.ListUp.PairFrame\_hsame\_injective}

\begin{theorem}[Pair frame empty head iff]
\label{thm:list-pair-frame-empty-head-iff}
For histories $u,t:\mathsf{BHist}$, the pair frame has the empty-head boundary
exactly when the head itself is empty:
$$
\hsame(\mathsf{PairFrame}(u,t),\Ezero(t))
\Longleftrightarrow
\hsame(u,\emp).
$$
\end{theorem}
\begin{proof}
Case-split the head $u$. In the empty case the pair-frame definition is
$\Ezero(t)$, so the displayed classifier is reflexive. In the zero-headed and
one-headed cases the pair frame begins with $\Eone$, contradicting
$\Eone/\Ezero$ no-confusion. Conversely, a classifier $\hsame(u,\emp)$
substitutes the empty head into the pair-frame definition.
\end{proof}
\leanchecked{BEDC.Derived.ListUp.PairFrame\_empty\_head\_iff}

\begin{definition}[Framed list endpoint]
\label{def:framed-list-endpoint}
For finite BEDC-history spines, define the framed endpoint by recursion on the
spine:
$$
\begin{array}{rcl}
\mathsf{FramedListEndpoint}(\mathsf{nil}) &=& \Ezero(\emp),\\
\mathsf{FramedListEndpoint}(\mathsf{cons}(a,xs)) &=&
\Eone(\mathsf{PairFrame}(a,\mathsf{FramedListEndpoint}(xs))).
\end{array}
$$
\end{definition}
\leandef{BEDC.Derived.ListUp.FramedListEndpoint}

\begin{theorem}[Framed list endpoint no-confusion and cons inversion]
\label{thm:framed-list-endpoint-no-confusion-cons-inversion}
The ten-field fit is: traditional target finite-list endpoint readback; BEDC
source $\mathsf{BHist}$ with $\hsame$; carrier framed finite spines; classifier
endpoint $\hsame$; stability obligations nil/cons separation and cons inversion;
ledger constructor frame; core theorem this block; exactness framed endpoint
exposes list shape; standard bridge finite spines; status
$\mathsf{ProofTarget}$. For all histories $a,b:\mathsf{BHist}$ and finite
spines $xs,ys$,
$$
\neg\hsame(\mathsf{FramedListEndpoint}(\mathsf{nil}),
\mathsf{FramedListEndpoint}(\mathsf{cons}(a,xs))),
$$
$$
\neg\hsame(\mathsf{FramedListEndpoint}(\mathsf{cons}(a,xs)),
\mathsf{FramedListEndpoint}(\mathsf{nil})),
$$
and
$$
\hsame(\mathsf{FramedListEndpoint}(\mathsf{cons}(a,xs)),
\mathsf{FramedListEndpoint}(\mathsf{cons}(b,ys)))
$$
implies both $\hsame(a,b)$ and
$\hsame(\mathsf{FramedListEndpoint}(xs),\mathsf{FramedListEndpoint}(ys))$.
\end{theorem}
\begin{proof}
The nil endpoint has outer constructor $\Ezero$, while every cons endpoint has
outer constructor $\Eone$, so the two mixed cases contradict BEDC history
no-confusion in the $\Ezero/\Eone$ and $\Eone/\Ezero$ directions. For the
cons-cons clause, invert the outer $\Eone$ equality and obtain equality of the
two displayed pair frames. Lemma~\ref{lem:list-pair-frame-injectivity} applies:
its proof is an induction on the first frame argument, with constructor
no-confusion separating the empty and tagged nonempty cases and with recursive
inversion exposing the first and second frame arguments. The two projections of
that lemma give $\hsame(a,b)$ and
$\hsame(\mathsf{FramedListEndpoint}(xs),\mathsf{FramedListEndpoint}(ys))$.
\end{proof}
\leanchecked{BEDC.Derived.ListUp.FramedListEndpoint\_no\_confusion\_cons\_inversion}

\begin{lemma}[Pair frame congruence]
\label{lem:list-pair-frame-congruence}
For histories $u,u',t,t':\mathsf{BHist}$, if $\hsame(u,u')$ and
$\hsame(t,t')$, then
$$
\hsame(\mathsf{PairFrame}(u,t),\mathsf{PairFrame}(u',t')).
$$
\end{lemma}
\begin{proof}
Proceed by induction on the displayed witness $\hsame(u,u')$. In the empty
case, the goal is the congruence of $\Ezero$ applied to $\hsame(t,t')$. In the
$\Ezero$ case, apply the induction hypothesis to the inner pair frame and then
apply congruence for the two visible constructors $\Ezero$ and $\Eone$. The
$\Eone$ case is identical except that the inner tag is $\Eone$. Thus the frame
recursion is stable in the first argument, and the terminal empty branch carries
the second-argument $\hsame$ witness.
\end{proof}
\leanchecked{BEDC.Derived.ListUp.PairFrame\_congruence}

\begin{theorem}[Framed list endpoint classifier exactness]
\label{thm:framed-list-endpoint-classifier-exactness}
For finite BEDC-history spines $xs,ys$,
$$
\hsame(\mathsf{FramedListEndpoint}(xs),\mathsf{FramedListEndpoint}(ys))
\Longleftrightarrow
\mathsf{ListClassifierSpec}(\hsame,xs,ys).
$$
\end{theorem}
\begin{proof}
For the left-to-right direction, argue by induction on the first spine and
case analysis on the second. The nil-nil case gives the nil constructor of
$\mathsf{ListClassifierSpec}$. The mixed nil-cons and cons-nil cases contradict
Theorem~\ref{thm:framed-list-endpoint-no-confusion-cons-inversion}. In the
cons-cons case, the same theorem exposes $\hsame$ between the heads and
$\hsame$ between the framed tail endpoints; apply the induction hypothesis to
the tail equality and then use the cons constructor of
$\mathsf{ListClassifierSpec}(\hsame,-,-)$.

For the right-to-left direction, induct on the displayed
$\mathsf{ListClassifierSpec}(\hsame,xs,ys)$ witness. The nil branch is
reflexivity of the framed nil endpoint. In the cons branch, the head witness and
the induction hypothesis give congruent arguments for
$\mathsf{PairFrame}$ by Lemma~\ref{lem:list-pair-frame-congruence}; one final
$\Eone$ congruence yields equality of the framed cons endpoints.
\end{proof}
\leanchecked{BEDC.Derived.ListUp.FramedListEndpoint\_classifier\_exactness}

\begin{theorem}[Framed list endpoint length preservation]
\label{thm:framed-list-endpoint-length-preservation}
The ten-field fit is: traditional target finite-list framed endpoint readback;
BEDC source $\mathsf{BHist}$; carrier finite spines; classifier endpoint
$\hsame$; stability endpoint exactness; ledger visible framed endpoint; core
theorem length preservation; exactness via $\mathsf{ListClassifierSpec}$;
standard bridge finite spines; status $\mathsf{ProofTarget}$. For finite
BEDC-history spines $xs,ys$, if
$$
\hsame(\mathsf{FramedListEndpoint}(xs),\mathsf{FramedListEndpoint}(ys))
$$
holds, then
$$
\ell(xs)=\ell(ys).
$$
\end{theorem}
\begin{proof}
Apply Theorem~\ref{thm:framed-list-endpoint-classifier-exactness} from left to
right to convert the displayed endpoint equality into
$\mathsf{ListClassifierSpec}(\hsame,xs,ys)$. Then apply
Theorem~\ref{thm:list-classifier-length-eq} to obtain equality of the finite
spine lengths.
\end{proof}
\leanchecked{BEDC.Derived.ListUp.FramedListEndpoint\_length\_preservation}

\begin{definition}[Framed list-spine representation]
\label{def:framed-list-spine-representation}
Let $A:\mathsf{BHist}\to\mathsf{Prop}$ be a BEDC history source carrier. Define
$\mathsf{FramedListSpineRep}_A(h,xs)$ to mean that every entry of the finite
spine $xs$ satisfies $A$ and that
$$
\hsame(h,\mathsf{FramedListEndpoint}(xs)).
$$
\end{definition}
\leandef{BEDC.Derived.ListUp.FramedListSpineRep}

\begin{theorem}[Framed list-spine representation hsame transport]
\label{thm:framed-list-spine-rep-hsame-transport}
Let $A:\mathsf{BHist}\to\mathsf{Prop}$ be a BEDC history source carrier. If
$\mathsf{FramedListSpineRep}_A(h,xs)$ holds and $\hsame(h,h')$ holds, then
$\mathsf{FramedListSpineRep}_A(h',xs)$ holds.
\end{theorem}
\begin{proof}
The entrywise source predicate over $xs$ is unchanged. The endpoint witness for
$h'$ is obtained by composing $\hsame(h',h)$, the symmetry of the displayed
$\hsame(h,h')$, with the stored witness
$\hsame(h,\mathsf{FramedListEndpoint}(xs))$.
\end{proof}
\leanchecked{BEDC.Derived.ListUp.FramedListSpineRep\_hsame\_transport}

\begin{theorem}[Framed list-spine cons endpoint injectivity]
\label{thm:list-spine-rep-cons-endpoint-injective}
Let $A:\mathsf{BHist}\to\mathsf{Prop}$ be a BEDC history source carrier. If
$\mathsf{FramedListSpineRep}_A(h,\mathsf{cons}(a,xs))$ and
$\mathsf{FramedListSpineRep}_A(k,\mathsf{cons}(b,ys))$ hold, and
$\hsame(h,k)$ holds, then the two framed cons endpoints are $\hsame$, and their
head payloads and tail endpoints are separately $\hsame$:
$$
\hsame(\Eone(\mathsf{PairFrame}(a,\mathsf{FramedListEndpoint}(xs))),
\Eone(\mathsf{PairFrame}(b,\mathsf{FramedListEndpoint}(ys))))
$$
and
$$
\hsame(a,b)\land
\hsame(\mathsf{FramedListEndpoint}(xs),\mathsf{FramedListEndpoint}(ys)).
$$
\end{theorem}
\begin{proof}
Compose the stored endpoint witness for the first representation, the displayed
$\hsame(h,k)$ witness, and the stored endpoint witness for the second
representation. This gives $\hsame$ between the two cons endpoints. Applying
\autoref{thm:framed-list-endpoint-no-confusion-cons-inversion} to that cons-cons
endpoint equality exposes both the head equality $\hsame(a,b)$ and the tail
endpoint equality.
\end{proof}
\leanchecked{BEDC.Derived.ListUp.FramedListSpineRep\_cons\_endpoint\_injective}

\begin{theorem}[Framed list-spine nil-cons no-confusion]
\label{thm:framed-list-spine-rep-nil-cons-no-confusion}
Let $A:\mathsf{BHist}\to\mathsf{Prop}$ be a BEDC history source carrier. No
single history can be represented both by the framed nil spine and by a framed
cons spine:
$$
\mathsf{FramedListSpineRep}_A(h,\mathsf{nil})\to
\mathsf{FramedListSpineRep}_A(h,\mathsf{cons}(a,xs))\to
\mathsf{False}.
$$
\end{theorem}
\begin{proof}
Unpack the two framed representations and compose their endpoint classifiers
through symmetry and transitivity of $\hsame$. This gives
$\hsame(\mathsf{FramedListEndpoint}(\mathsf{nil}),
\mathsf{FramedListEndpoint}(\mathsf{cons}(a,xs)))$, which contradicts the
framed endpoint nil-cons no-confusion theorem.
\end{proof}
\leanchecked{BEDC.Derived.ListUp.FramedListSpineRep\_nil\_cons\_no\_confusion}
